\begin{theorem}[DP31: the scalar higher-height correction is integral]\label{DP31}
Let $h>2$, let $k$ contain $\F_q$, and now take $R=W(k)$ and $K=R[1/p]$. Under (H2)--(H4),
\begin{equation}\label{eq:scalar-gauge}
 B(S,T):=\exp\beta(L_h(S),L_h(T))\in1+ST\,R[[S,T]].
\end{equation}
Consequently the scalar function of DP29 satisfies the stronger estimate
\begin{equation}\label{eq:scalar-full-slope}
 v_p(\kappa![\mathbf t^\kappa]\Theta)
 \ge c_h|\kappa|-\frac{\sum_i s_p(\kappa_i)}{p-1}.
\end{equation}
The slope $c_h$ is optimal within this class, by \eqref{eq:sharp-corrected-scalar}. This theorem concerns two scalar copies of the height-$h$ group, not the genuinely dual pair.
\end{theorem}
\begin{proof}
Let $G$ have law $F_h$, and write $R_n=R/p^nR$. We use the established formal biextension theorem: Zink \cite[\S4.4, Proposition 121 and (243), pp.~241, 246]{Zink} gives an injection
\[
 \operatorname{Biext}^1(G_n\times G_n,\widehat{\Gm})
 \hookrightarrow\Hom(G_n,G_n^t).
\]
The formal groups here come from displays; the existence/equivalence applies over the Artinian local rings $R_n$ by \cite[Theorem 9, or Theorem 103]{Zink}. The dual display gives the Cartier-dual formal group. Both $G$ and its Cartier dual are connected for $h>1$, with special-fibre slopes $1/h$ and $(h-1)/h$ respectively. The hypotheses of the injection are therefore satisfied; we do not need the additional surjectivity in Zink's Theorem 122.

First, the integral units $C_X=H_X$ and $C_Y=H_Y$ define a biextension without presupposing \eqref{eq:scalar-gauge}. Over $K$ they are the two multiplicative coboundaries of $B$. Therefore they satisfy symmetry, associativity in each variable, and the interchange identity. These identities hold over $R$, since $R\hookrightarrow K$. On the trivial formal multiplicative torsor over $G\times G$, define the first partial product by
\[
 (S,T,u)\star_X(S',T,u')
  =(F_h(S,S'),T,uu'C_X(S,S',T)),
\]
and define the second analogously with $C_Y$. The identities just stated make a biextension. A normalized splitting is exactly a unit $I$ with
\begin{equation}\label{eq:gauge-split}
 \frac{I(F_h(S,S'),T)}{I(S,T)I(S',T)}=C_X,\qquad
 \frac{I(S,F_h(T,T'))}{I(S,T)I(S,T')}=C_Y,
 \qquad I(S,0)=I(0,T)=1.
\end{equation}

Second, $\Hom(G_n,G_n^t)=0$. Over $k$, distinct slopes give zero rational Hom, hence zero integral Hom. The integral Hom embeds in its rationalization: a homomorphism killed by $p^a$ is zero since multiplication by $p^a$ on a $p$-divisible group is surjective as a sheaf. To lift the vanishing, let $f:G_n\to G_n^t$ reduce to zero, and suppose its formal-coordinate components are in $p^jR_n[[T]]$, $1\le j<n$. Write $f=p^jg$ componentwise. The target $p$-series has linear part $p$ times the identity and higher terms of order at least two, so
\[
 f([p]_G(T))=[p]_{G^t}(f(T))\equiv0\pmod {p^{j+1}}.
\]
After division by $p^j$ and reduction, each component satisfies $\bar g([p]_{G_k}(T))=0$. Substitution of the nonzero series $[p]_{G_k}(T)$, whose initial degree is $q$, is injective on $k[[T]]$. Thus $f$ is divisible by $p^{j+1}$, and induction proves $f=0$.

The biextension injection now yields a splitting $I_n\in1+ST\,R_n[[S,T]]$ at every $n$. We verify compatibility rather than assume that independently chosen splittings lift. Two normalized splittings differ by a bihomomorphism $J$, which obeys
\[
 J([p]_G(S),T)=J(S,T)^p.
\]
Modulo $p$, let $b>0$ be the smallest $T$-degree in $J-1$, if it is nonzero. Its coefficient on the left remains nonzero under substitution of $[p]_{G_k}(S)$; on the right every positive $T$-degree is at least $pb>b$. This contradiction proves $J=1$ modulo $p$. If $J-1=p^jV$ for $j\ge1$, the right side minus one is divisible by $p^{j+1}$. Dividing and reducing gives $\bar V([p]_{G_k}(S),T)=0$, and injectivity of this substitution on $k[[S,T]]$ proves the next divisibility. Hence $J=1$ over every $R_n$.

Thus the normalized splittings are unique and automatically compatible. Their coefficientwise inverse limit is $I\in1+ST\,R[[S,T]]$ satisfying \eqref{eq:gauge-split}. In additive coordinates,
\[
 A(X,Y)=\beta(X,Y)-\log I(E_h(X),E_h(Y))
\]
is separately additive, by equality of the two coboundaries. In characteristic zero, interior binomial coefficients force a separately additive series to be $aXY$. The logarithm and exponential are Teichm\"uller-equivariant. Consequently the assumed torus symmetry of $\beta$ makes $I_n(\zeta S,\zeta^{-p}T)$ another normalized splitting. Uniqueness makes $I$, and hence $A$, invariant. But $XY$ has character $\zeta^{1-p}\ne1$, so $a=0$. This proves $B=I$ and \eqref{eq:scalar-gauge}.

Finally, write the correction factor in $\Theta$ as
\[
 \prod_j B(E_h(\ell_j),E_h(\epsilon\ell_j))^{\sigma/2}.
\]
Each binomial power is an integral series because $p$ is odd. The whole expression is an integral function of finitely many inverse-logarithm coordinates of weight at least $\rho_h$. Therefore $v_p([\mathbf t^\kappa]\Theta)\ge-\rho_h|\kappa|$. Legendre's factorial formula gives \eqref{eq:scalar-full-slope}. The explicit nonzero example in DP29 attains the slope along $q$-power orders.
\end{proof}
